% !PRINTED 626
\weberpage{626}
b) If, conversely, $\gamma$ is divisible by $q-1$, it follows from (8) that $p-1$ is divisible by $q$. If, then, we assume that $p-1$ is not divisible by $q$, then $e$ is also not divisible by $q-1$, and $\xi$ runs through the sequence of numbers (9).

If we form the sum of these numbers $\xi$, it follows that:
\[ (g-1) \Sigma \xi \equiv g^e - 1 \, (\mathrm{mod}\ m), \]
and this number is divisible or not divisible by $q$, according as $e$ is divisible or not divisible by $q-1$. Since, moreover, $g-1$ is not divisible by $q$, $\Sigma\xi$ is divisible or not divisible by $q$, according as $p-1$ is divisible or not divisible by $q$.

2. We turn to the case in which $m = 2^\varkappa$ is a power of 2, and here we have again to distinguish whether $p-1$ is divisible by 4 or only by 2.

a) If $p-1$ is divisible by 4, then (\S.\,15):
\[ p \equiv 5^\beta \, (\mathrm{mod}\ m), \]
and $f$ is the least positive number that satisfies the condition
\[ \beta f \equiv 0 \, [\mathrm{mod}\ \tfrac{1}{2}\varphi(m)]. \]
If $f=1$, that is, if $\beta$ is divisible by $\tfrac{1}{2}\varphi(m) = 2^{\varkappa-2}$, then $e = \varphi(m)$, and $\xi$ runs through the whole group $\mathfrak{N}$, i.\,e. all odd numbers between 0 and $m$. This is the case in which $p-1$ is divisible by $m$.

If $2^{\varkappa_1-2}$ is the highest power of 2 that divides $\beta$, and $2 \leqq \varkappa_1 < \varkappa$, then, if we set $m_1 = 2^{\varkappa_1}$ and $m = m_1 m_2$,
\[ f = 2^{\varkappa - \varkappa_1}, \quad e = 2^{\varkappa_1 - 1} = \varphi(m_1), \]
and, if $s$ is an odd number, we have
\[ p - 1 \equiv 5^{s\tfrac{1}{2}\varphi(m_1)} - 1 \ldots \, (\mathrm{mod}\ m). \]
% [sic] Source prints three low dots after -1 in this congruence; no continuation term is specified.
By \S.\,15, $5^{\tfrac{1}{2}\varphi(m_1)} - 1$ is divisible by $m_1$, but by no higher power of 2. Since, moreover, one has by the same theorems
\[ 5^{s\tfrac{1}{2}\varphi(m_1)} \equiv 5^{\tfrac{1}{2}\varphi(m_1)} \, (\mathrm{mod}\ 2m_1) \]
it follows that $m_1$ is the greatest common divisor of $p-1$ and $m$.

If, then, as above, we call in the field $\Omega_{m_1}$, it results that, when $m_1 \geqq 4$ is the greatest common divisor of $p-1$ and $m$, $\xi$ runs through a complete system of odd residues of $m_1$, and that $\mathfrak{p}_\xi$ are at the same time the prime factors of $p$ in the field $\Omega_{m_1}$.
